%%%%%%%%%%%%%%%%%%%%%%%%%%%%%%%%%%%%%%%%%%%%%%%%%%%%%%%
% CHAPTER: Preliminaries
%%%%%%%%%%%%%%%%%%%%%%%%%%%%%%%%%%%%%%%%%%%%%%%%%%%%%%%
\chapter{Preliminaries}
\label{chap:pre}

% —————————————————————————————————————————————————————
% 
% —————————————————————————————————————————————————————

%%%%%%%%%%%%%%%%%%%%%%%%%%%%%%%%%%%%%%%%%%%%%%%%%%%%%%%
% Mathematical Preliminaries
%%%%%%%%%%%%%%%%%%%%%%%%%%%%%%%%%%%%%%%%%%%%%%%%%%%%%%%
\section{Mathematical Preliminaries}

We review some mathematical preliminaries that will be used throughout this proposal.

\begin{definition}
  \label{PRE:def:C11}
  A class $\mathcal{C}^{1,1}$ of functions is defined as the set of functions $\mv{f}:\mathbb{R}^n\to\mathbb{R}^m$ that are continuously differentiable and whose gradients are Lipschitz continuous.
\end{definition}

\begin{lem}[{\cite[Lemma 3.1]{Khalil:2002aa}}]
  \label{PRE:lem:con:lip}
  Consider a function $\mv{f}:\mathbb{R}^n\to\mathbb{R}^m$ that is continuously differentiable on a compact convex set $\mathcal{X}\subset\mathbb{R}^n$. 
  Then, $\mv{f}$ is Lipschitz continuous on $\mathcal{X}$, implying
  \begin{equation}
    \exists
    L\in\R_{>0}
    :
    \norm{\mv{f}(\mv{x}_1)-\mv{f}(\mv{x}_2)}\le L\norm{\mv{x}_1-\mv{x}_2}
    ,
    \quad\forall \mv{x}_1,\mv{x}_2\in\mathcal{X}
    .
  \end{equation}
\end{lem}

\begin{proof}

  Since $\mv{f}$ is continuously differentiable on a compact set $\mathcal{X}$, the gradient of $\mv{f}$ is bounded on $\mathcal{X}$.
  Therefore, there exists a constant $L>0$ such that $\norm{\nabla \mv{f}(\mv{x})}\le L$ for all $\mv{x}\in\mathcal{X}$, where $\nabla \mv{f}(\mv{x}):=\ppfrac{\mv{f}(\mv{x})}{\mv{x}}$ denotes the Jacobian matrix of $\mv{f}$ at $\mv{x}$.
  By the mean value theorem, for any $\mv{x}_1,\mv{x}_2\in\mathcal{X}$, there exists a point $\mv{c}$ on the line segment connecting $\mv{x}_1$ and $\mv{x}_2$ such that
  \begin{equation}
    \mv{f}(\mv{x}_1)-\mv{f}(\mv{x}_2) = \nabla \mv{f}(\mv{c})^\top (\mv{x}_1-\mv{x}_2)
    .
  \end{equation}
  Taking the norm of both sides and applying the bound on the gradient, we have
  \begin{equation}
    \norm{\mv{f}(\mv{x}_1)-\mv{f}(\mv{x}_2)} 
    \le 
    \norm{\nabla \mv{f}(\mv{c})}
    \norm{\mv{x}_1-\mv{x}_2} 
    \le 
    L\norm{\mv{x}_1-\mv{x}_2}
    ,
  \end{equation}
  which concludes the proof.

\end{proof}

\begin{lem}
  \label{PRE:lem:lib:grad:bound}
  Consider a function $\mv{f}:\mathbb{R}^n\to\mathbb{R}^m$ that is $L$-Lipschitz continuous on a compact set $\mathcal{X}\subset\mathbb{R}^n$. 
  Then, the gradient of $\mv{f}$ is bounded on $\mathcal{X}$, implying
  \begin{equation}
    \exists
    M\in\R_{>0}
    :
    \norm{
      \nabla \mv{f}(\mv{x})
    }\le M
    ,
    \quad\forall \mv{x}\in\mathcal{X}
    ,
  \end{equation}
  where $\nabla \mv{f}(\mv{x}):=\ppfrac{\mv{f}(\mv{x})}{\mv{x}}$ denotes the Jacobian matrix of $\mv{f}$ at $\mv{x}$.
\end{lem}

\begin{proof}
  \cred{To be added.}
\end{proof}

\begin{lem}
  \label{PRE:lem:lib:ineq}
  Consider a function $\mv{f}:\mathbb{R}^n\to\mathbb{R}^m$ that is $L$-Lipschitz continuous on a compact set $\mathcal{X}\subset\mathbb{R}^n$. 
  Then, the following inequality holds for all $\mv{x}_1,\mv{x}_2\in\mathcal{X}$:
  \begin{equation}
    \mv{f}(\mv{x}_1)
    \le 
    \mv{f}(\mv{x}_2)
    +
    \nabla\mv{f}(\mv{x}_2)^\top
    (\mv{x}_1-\mv{x}_2)
    +
    \frac{L}{2}
    \norm{\mv{x}_1-\mv{x}_2}^2
    .
  \end{equation}
\end{lem}

\begin{proof}

  A $L$-Lipschitz continuous function $\mv{f}$ satisfies the following inequality for all $\mv{x}_1,\mv{x}_2\in\mathcal{X}$:
  \begin{equation}
    \label{PRE:eq:lib:ineq:1}
    \norm{\mv{f}(\mv{x}_1) - \mv{f}(\mv{x}_2)}
    \le
    L\norm{\mv{x}_1-\mv{x}_2}
    .
  \end{equation}

  Additionally, we have
  \begin{equation}
    \label{PRE:eq:lib:ineq:2}
    f(\mv{x}_2)
    =
    f(\mv{x}_1)
    +
    \nabla f(\mv{x}_1)^\top
    \left(
      \mv{x}_2 - \mv{x}_1
    \right)
    +
    \int_{0}^{1}
    \big(
      \nabla f(\mv{x}_1 + \tau(\mv{x}_2-\mv{x}_1))
      -
      \nabla f(\mv{x}_1)
    \big)^\top
    \left(
      \mv{x}_2 - \mv{x}_1
    \right)
    \,\der\tau
  \end{equation}
  Using the Cauchy-Schwarz inequality and \eqref{PRE:eq:lib:ineq:1}, we have
  \begin{equation}
    \label{PRE:eq:lib:ineq:3}
    \begin{aligned}
      \int_{0}^{1}
      \big(
        \nabla f(\mv{x}_1 + \tau(\mv{x}_2-\mv{x}_1))
        -
        \nabla f(\mv{x}_1)
      \big)^\top
      \left(
        \mv{x}_2 - \mv{x}_1
      \right)
      \,\der\tau
      \le
      &
      \int_{0}^{1}
      \norm{
        \nabla f(\mv{x}_1 + \tau(\mv{x}_2-\mv{x}_1))
        -
        \nabla f(\mv{x}_1)
      }
      \norm{
        \mv{x}_2 - \mv{x}_1
      }
      \,\der\tau
      \\
      \le
      &
      \int_{0}^{1}
      L\tau\norm{\mv{x}_2-\mv{x}_1}^2
      \,\der\tau
      =
      \frac{L}{2}
      \norm{\mv{x}_2-\mv{x}_1}^2
      .
    \end{aligned}
  \end{equation}
  Finally, by substituting \eqref{PRE:eq:lib:ineq:2} into \eqref{PRE:eq:lib:ineq:1}, we can conclude the proof by presenting the desired result.

\end{proof}

\begin{lem}
  \label{PRE:lem:lib:grad:bound:PL}
  Consider a function $\mv{f}:\mathbb{R}^n\to\mathbb{R}^m$ that is $L$-Lipschitz continuous.
  In addition, assume that $\mv{f}$ satisfies $\mu$-Polyak-Łojasiewicz (PL) condition on a compact set $\mathcal{X}\subset\mathbb{R}^n$, implying
  \begin{equation}
    \exists \mu\in\R_{>0}
    :
    \norm{\nabla \mv{f}(\mv{x})}^2
    \ge
    2\mu
    \left(
      \mv{f}(\mv{x}) - \mv{f}^*
    \right)
    ,
    \quad\forall \mv{x}\in\mathcal{X}
    ,
  \end{equation}
  where $\mv{f}^*:=\min_{\mv{x}\in\mathcal{X}}\mv{f}(\mv{x})$.
  Then, the following inequality holds:
  \begin{equation}
    L \ge \mu
    . 
  \end{equation}
\end{lem}

\begin{proof}

  From Lemma~\ref{PRE:lem:lib:ineq}, we have
  \begin{equation}
    \mv{f}(\mv{x}_1)
    \le 
    \mv{f}(\mv{x}_2)
    +
    \nabla\mv{f}(\mv{x}_2)^\top
    (\mv{x}_1-\mv{x}_2)
    +
    \frac{L}{2}
    \norm{\mv{x}_1-\mv{x}_2}^2
    ,
    \quad\forall \mv{x}_1,\mv{x}_2\in\mathcal{X}
    .
  \end{equation}
  By substituting $\mv{x}_1=\mv{x}_2-\frac{1}{L}\nabla\mv{f}(\mv{x}_2)$ into the above inequality, we have
  \begin{equation}
    \mv{f}(\mv{x}_1)
    \le
    \mv{f}(\mv{x}_2)
    -
    \frac{1}{2L}
    \norm{\nabla\mv{f}(\mv{x}_2)}^2
    ,
    \quad\forall \mv{x}_1,\mv{x}_2\in\mathcal{X}
    .
  \end{equation}

  By rearranging the above inequality, we have
  \begin{equation}
    \norm{\nabla\mv{f}(\mv{x}_2)}^2
    \le
    2L
    \left(
      \mv{f}(\mv{x}_2) - \mv{f}(\mv{x}_1)
    \right)
    ,
    \quad\forall \mv{x}_1,\mv{x}_2\in\mathcal{X}
    .
  \end{equation}
  Since $\mv{f}$ satisfies the $\mu$-PL condition, we have
  \begin{equation}
    2\mu
    \left(
      \mv{f}(\mv{x}_2) - \mv{f}^*
    \right)
    \le
    \norm{\nabla\mv{f}(\mv{x}_2)}^2
    \le
    2L
    \left(
      \mv{f}(\mv{x}_2) - \mv{f}(\mv{x}_1)
    \right)
    ,
    \quad\forall \mv{x}_2\in\mathcal{X}
    .
  \end{equation}
  Since $\mv{f}(\mv{x}_1)\ge \mv{f}^*$, we have
  \begin{equation}
    2\mu
    \left(
      \mv{f}(\mv{x}_2) - \mv{f}^*
    \right)
    \le
    2L
    \left(
      \mv{f}(\mv{x}_2) - \mv{f}^*
    \right)
    ,
    \quad\forall \mv{x}_2\in\mathcal{X}
    ,
  \end{equation}
  which implies that $L\ge \mu$.

\end{proof}

%%%%%%%%%%%%%%%%%%%%%%%%%%%%%%%%%%%%%%%%%%%%%%%%%%%%%%%
% Neural Network Preliminaries
%%%%%%%%%%%%%%%%%%%%%%%%%%%%%%%%%%%%%%%%%%%%%%%%%%%%%%%
\section{Neural Network Preliminaries}
\label{PRE:sec:NN}

The unconstrained deep neural networks (DNNs) $\NN_u(\NNin;\wth)$ is parameterized by the weight vector $\wth\in\R^{\Xi}$ and the input vector $\NNin\in\R^{l_0+1}$.
Specifically, the unconstrained DNN $\NN_u:=\NN_u(\NNin;\wth)$ is recursively defined as follows:
\begin{equation}
  \label{PRE:eq:DNN}
  \NN_l
  =
  \begin{cases}
    \wV_l^\top\act_l
    ,
    & 
    l \in \{1,\ldots,L\}
    \\
    \wV_0^\top\NNin
    ,
    &
    l = 0
    ,
  \end{cases}
\end{equation}
where $\wV_l\in\R^{(\ell_{l}+1)\times \ell_{l+1}}$ is the weight matrix, and $\act_l:\R^{\ell_{l}}\to\R^{\ell_{l}+1}$ is the activation vector of the $l$-th layer.
To incorporate the bias term, the last term of the input vector $\NNin$ is set to $1$, while the remaining terms are designed to include all arguments of the function to be approximated.
Similarly, the activation vector $\act_l(\mv{z})$ is defined as 
$
  \act_i(\mv{z}) 
  = 
  \begin{pmatrix}
  \phi_1(z_1) & \phi_2(z_2) & \cdots & \phi_{l_i}(z_{l_i}) & 1
  ,
  \end{pmatrix}^\top
$
where $\phi_j:\R\to\R$ is the activation function of the $j$-th neuron in the $l$-th layer, for some $\mv{z}:=(z_1,\ldots,z_{l_i})^\top\in\R^{l_i}$.

\begin{lem}[Universal approximation property of DNNs]
  \label{PRE:lem:NN:UAP}
  Let $\Omega_{nn}\subset\R^n$ be a compact set, and let $\mv{f}:\Omega_{nn}\to\R^m$ be a continuous function. 
  Then, for any $\epsilon>0$, there exists a neural network $\NN_u(\NNin;\wth)$ with a finite number of neurons such that
  \begin{equation}
    \sup_{\NNin\in\Omega_{nn}}
    \norm{
      \mv{f}(\NNin) - \NN_u(\NNin;\wth)
    }
    <
    \epsilon
    .
  \end{equation}
\end{lem}

\begin{proof}
  \cred{To be added.}
\end{proof}

% Lipschitz continuity of unconstrained DNNs
\begin{lem}[Lipschitz continuity of a neural network with respect to its parameters]
  \label{PRE:lem:NN:param:Lipschitz}
  Consider the activation functions $\phi_l,\ \forall l\in\{1,\ldots,L\}$, of the unconstrained DNN $\NN_u(\NNin;\wth)$ in \eqref{PRE:eq:DNN} that are Lipschitz continuous with respect to their arguments. 
  Then, the unconstrained DNN $\NN_u(\NNin;\wth)$ is Lipschitz continuous with respect to its weight vector $\wth$ on a compact set $\mathcal{W}\subset\R^{\Xi}$, implying
  \begin{equation}
    \exists
    L_{\NN_u}\in\R_{>0}
    :
    \norm{
      \NN_u(\NNin;\wth_1) - \NN_u(\NNin;\wth_2)
    }
    \le
    L_{\NN_u}
    \norm{\wth_1-\wth_2}
    ,
    \quad\forall \wth_1,\wth_2\in\mathcal{W}
    .
  \end{equation}
\end{lem}

\begin{proof}

  The gradient of $\NN_u$ with respect to $\wth$ can be obtained using the chain rule as follows:
  \begin{equation}
    \ppfrac{\NN}{\wth}=
    \begin{bmatrix}
      \ppdfrac{\NN}{\wth_k}&
      \ppdfrac{\NN}{\wth_{k-1}}&
  \cdots &
      \ppdfrac{\NN}{\wth_0}
    \end{bmatrix}
    \in\R^{l_{k+1} \times \Xi}
    ,
    \label{PRE:eq:NN:jacobian}
  \end{equation}
  where
  \begin{equation}
    \ppfrac{\NN}{\wth_i} = 
    \begin{cases}
      % (
        \mm I_{l_{k+1}}
        \otimes 
        \act_{k}^\top  
      % )
      , 
      &
      i=k
      ,
      \\
      \wV_k^\top\act_{k}' 
      (
        \mm I_{l_{k}}
        \otimes  
        \act_{k-1}^\top  
      )
      , 
      & 
      i=k-1
      ,
      \\
      \vdots 
      &
      \vdots 
      \\
      \wV_k^\top\act'_{k} 
      \cdots 
      \wV_1^\top\act_1' 
      (
        \mm I_{l_1}
        \otimes 
        \NNin^\top  
      )
      , 
      &
      i = 0
      ,
    \end{cases}
  \end{equation}
  and $\act_i:= \act_i(\NN_{i-1})$ and $\act_i':= \ddtfrac{\act_i}{\NN_{i-1}}$, for all $i\in\mathcal{K}\setminus\{0\}$.

  \cred{To be added.}

\end{proof}

% % Lipschitz continuity of the unconstrained DNN
% \begin{lem}[Lipschitz continuity of a neural network with respect to its parameters]
%   \label{PRE:lem:NN:param:Lipschitz}
%   Consider the unconstrained DNN
%   $\NN_u(\NNin;\wth)$ defined in \eqref{PRE:eq:DNN} for a fixed input
%   $\NNin$, where
%   \begin{equation}
%     \wth
%     =
%     \operatorname{col}
%     \left(
%       \wth_0,\wth_1,\ldots,\wth_L
%     \right)
%     \in\R^\Xi.
%   \end{equation}
%   Suppose that each activation function
%   $\phi_l$, $l\in\{1,\ldots,L\}$, is
%   $L_{\phi_l}$-Lipschitz continuous and that
%   $\wth$ belongs to a compact set
%   $\mathcal{W}\subset\R^\Xi$.
%   Then, there exists a constant
%   $L_{\NN_u}\in\R_{>0}$ such that
%   \begin{equation}
%     \label{PRE:eq:NN:param:Lipschitz}
%     \norm{
%       \NN_u(\NNin;\wth_1)
%       -
%       \NN_u(\NNin;\wth_2)
%     }
%     \leq
%     L_{\NN_u}
%     \norm{\wth_1-\wth_2},
%     \qquad
%     \forall\wth_1,\wth_2\in\mathcal{W}.
%   \end{equation}
%   Moreover, $\NN_u$ is differentiable with respect to $\wth$
%   almost everywhere and satisfies
%   \begin{equation}
%     \norm{A}
%     \leq
%     L_{\NN_u},
%     \qquad
%     \forall
%     A\in
%     \partial_C\NN_u(\NNin;\wth),
%     \quad
%     \forall\wth\in\mathcal{W},
%   \end{equation}
%   where $\partial_C$ denotes the Clarke generalized Jacobian.
% \end{lem}

% \begin{proof}
%   Define the compact convex set
%   \begin{equation}
%     \widetilde{\mathcal{W}}
%     :=
%     \operatorname{co}(\mathcal{W}),
%   \end{equation}
%   where $\operatorname{co}(\mathcal{W})$ denotes the convex hull of
%   $\mathcal{W}$.
%   Since the weight matrices are components of $\wth$, their norms are
%   uniformly bounded on $\widetilde{\mathcal{W}}$.
%   Thus, there exist constants
%   $\overline{V}_i\in\R_{>0}$ such that
%   \begin{equation}
%     \norm{\wV_i}
%     \leq
%     \overline{V}_i,
%     \qquad
%     \forall\wth\in\widetilde{\mathcal{W}}.
%   \end{equation}

%   Since the affine mappings and activation functions are continuous,
%   all layer outputs are continuous with respect to $\wth$.
%   Their continuity on the compact set
%   $\widetilde{\mathcal{W}}$ implies that there exist constants
%   $\overline{a}_i\in\R_{>0}$ satisfying
%   \begin{equation}
%     \norm{\act_i}
%     \leq
%     \overline{a}_i,
%     \qquad
%     \forall\wth\in\widetilde{\mathcal{W}},
%   \end{equation}
%   where
%   \begin{equation}
%     \act_0:=\NNin,
%     \qquad
%     \overline{a}_0:=\norm{\NNin}.
%   \end{equation}

%   Since $\phi_i$ is $L_{\phi_i}$-Lipschitz continuous,
%   Rademacher's theorem implies that it is differentiable almost
%   everywhere and
%   \begin{equation}
%     \norm{\act_i'}
%     \leq
%     L_{\phi_i}
%   \end{equation}
%   wherever the derivative exists.
%   The same bound holds for every element of its Clarke generalized
%   Jacobian at a nondifferentiability point.

%   At a differentiability point, the Jacobian of the DNN with respect
%   to the weight block $\wth_i$ is
%   \begin{equation}
%     \ppfrac{\NN_u}{\wth_i}
%     =
%     \begin{cases}
%       I_{l_{L+1}}\otimes\act_L^\top,
%       & i=L,
%       \\[1mm]
%       \wV_L^\top\act_L'
%       \cdots
%       \wV_{i+1}^\top\act_{i+1}'
%       \left(
%         I_{l_{i+1}}\otimes\act_i^\top
%       \right),
%       & i\in\{0,\ldots,L-1\}.
%     \end{cases}
%   \end{equation}
%   Therefore, using the submultiplicativity of the induced matrix norm
%   and
%   \begin{equation}
%     \norm{I_p\otimes z^\top}
%     =
%     \norm{z},
%   \end{equation}
%   we obtain
%   \begin{equation}
%     \norm{
%       \ppfrac{\NN_u}{\wth_i}
%     }
%     \leq
%     \overline{a}_i
%     \prod_{j=i+1}^{L}
%     \left(
%       \overline{V}_jL_{\phi_j}
%     \right)
%     =:
%     G_i,
%     \qquad
%     i\in\{0,\ldots,L\},
%   \end{equation}
%   where the empty product for $i=L$ is defined as one.

%   The complete Jacobian is given by
%   \begin{equation}
%     \ppfrac{\NN_u}{\wth}
%     =
%     \begin{bmatrix}
%       \ppdfrac{\NN_u}{\wth_L}
%       &
%       \ppdfrac{\NN_u}{\wth_{L-1}}
%       &
%       \cdots
%       &
%       \ppdfrac{\NN_u}{\wth_0}
%     \end{bmatrix}.
%   \end{equation}
%   Hence,
%   \begin{equation}
%     \norm{
%       \ppfrac{\NN_u}{\wth}
%     }
%     \leq
%     \left(
%       \sum_{i=0}^{L}G_i^2
%     \right)^{1/2}
%     =:
%     L_{\NN_u}.
%   \end{equation}
%   By the Clarke chain rule, the same bound holds at
%   nondifferentiability points:
%   \begin{equation}
%     \norm{A}
%     \leq
%     L_{\NN_u},
%     \qquad
%     \forall
%     A\in\partial_C\NN_u(\NNin;\wth).
%   \end{equation}

%   For arbitrary $\wth_1,\wth_2\in\mathcal{W}$, define
%   \begin{equation}
%     \gamma(s)
%     :=
%     \wth_2+s(\wth_1-\wth_2),
%     \qquad
%     s\in[0,1].
%   \end{equation}
%   Since
%   $\gamma(s)\in\widetilde{\mathcal{W}}$ and
%   $\NN_u(\NNin;\gamma(s))$ is absolutely continuous, the generalized
%   chain rule gives, for almost every $s\in[0,1]$,
%   \begin{equation}
%     \frac{d}{ds}
%     \NN_u(\NNin;\gamma(s))
%     =
%     A(s)(\wth_1-\wth_2)
%   \end{equation}
%   for some
%   $A(s)\in\partial_C\NN_u(\NNin;\gamma(s))$.
%   Consequently,
%   \begin{align}
%     &
%     \norm{
%       \NN_u(\NNin;\wth_1)
%       -
%       \NN_u(\NNin;\wth_2)
%     }
%     \nonumber\\
%     &\quad\leq
%     \int_0^1
%     \norm{A(s)}
%     \norm{\wth_1-\wth_2}
%     \,ds
%     \nonumber\\
%     &\quad\leq
%     L_{\NN_u}
%     \norm{\wth_1-\wth_2},
%   \end{align}
%   which proves \eqref{PRE:eq:NN:param:Lipschitz}.
% \end{proof}

% ```latex
% \begin{lem}[Lipschitz continuity of a DNN with respect to its parameters]
%   \label{PRE:lem:NN:param:Lipschitz}
%   Consider the feedforward neural network
%   \begin{equation}
%     \begin{aligned}
%       \boldsymbol{h}_0
%       &=
%       \boldsymbol{x},
%       \\
%       \boldsymbol{h}_l
%       &=
%       \phi_l
%       \left(
%         W_l\boldsymbol{h}_{l-1}+\boldsymbol{b}_l
%       \right),
%       \qquad l\in\{1,\ldots,L\},
%       \\
%       \Phi(\boldsymbol{x};\boldsymbol{\theta})
%       &=
%       W_{L+1}\boldsymbol{h}_L+\boldsymbol{b}_{L+1},
%     \end{aligned}
%   \end{equation}
%   where $\boldsymbol{\theta}$ collects all weights and biases.
%   Suppose that each activation function $\phi_l$ is Lipschitz continuous with
%   constant $L_{\phi_l}$ and that
%   $\boldsymbol{\theta}$ belongs to a compact set
%   $\Theta\subset\mathbb{R}^{N_\theta}$.
%   Then, for every fixed input $\boldsymbol{x}$, there exists a constant
%   $L_\Phi(\boldsymbol{x})>0$ such that
%   \begin{equation}
%     \left\|
%       \Phi(\boldsymbol{x};\boldsymbol{\theta}_1)
%       -
%       \Phi(\boldsymbol{x};\boldsymbol{\theta}_2)
%     \right\|
%     \leq
%     L_\Phi(\boldsymbol{x})
%     \left\|
%       \boldsymbol{\theta}_1-\boldsymbol{\theta}_2
%     \right\|,
%     \qquad
%     \forall\boldsymbol{\theta}_1,\boldsymbol{\theta}_2\in\Theta.
%   \end{equation}
%   In particular, the result holds for the ReLU activation function with
%   $L_{\phi_l}=1$.
% \end{lem}

% \begin{proof}
%   Since $\Theta$ is compact, the network parameters are uniformly bounded.
%   Thus, there exist constants $\overline{W}_l,\overline{b}_l>0$ such that
%   \begin{equation}
%     \|W_l\|\leq\overline{W}_l,
%     \qquad
%     \|\boldsymbol{b}_l\|\leq\overline{b}_l,
%     \qquad
%     \forall\boldsymbol{\theta}\in\Theta.
%   \end{equation}
%   For the fixed input $\boldsymbol{x}$, define
%   $\overline{h}_0:=\|\boldsymbol{x}\|$ and recursively select
%   $\overline{h}_l$ such that
%   \begin{equation}
%     \overline{h}_l
%     :=
%     \|\phi_l(\boldsymbol{0})\|
%     +
%     L_{\phi_l}
%     \left(
%       \overline{W}_l\overline{h}_{l-1}
%       +
%       \overline{b}_l
%     \right).
%   \end{equation}
%   The Lipschitz continuity of $\phi_l$ then gives
%   \begin{equation}
%     \|\boldsymbol{h}_l\|
%     \leq
%     \overline{h}_l,
%     \qquad
%     \forall\boldsymbol{\theta}\in\Theta.
%   \end{equation}

%   Let $\boldsymbol{h}_l^i$ denote the output of layer $l$ corresponding to
%   $\boldsymbol{\theta}_i$, where $i\in\{1,2\}$. Since the input is fixed,
%   \begin{equation}
%     \|\boldsymbol{h}_0^1-\boldsymbol{h}_0^2\|=0.
%   \end{equation}
%   For each hidden layer, the Lipschitz continuity of $\phi_l$ yields
%   \begin{align}
%     \|\boldsymbol{h}_l^1-\boldsymbol{h}_l^2\|
%     \leq
%     L_{\phi_l}
%     \Big(
%       &
%       \overline{W}_l
%       \|\boldsymbol{h}_{l-1}^1-\boldsymbol{h}_{l-1}^2\|
%       \nonumber\\
%       &+
%       (\overline{h}_{l-1}+1)
%       \|\boldsymbol{\theta}_1-\boldsymbol{\theta}_2\|
%     \Big).
%   \end{align}
%   Define $G_0:=0$ and
%   \begin{equation}
%     G_l
%     :=
%     L_{\phi_l}
%     \left(
%       \overline{W}_lG_{l-1}
%       +
%       \overline{h}_{l-1}
%       +
%       1
%     \right).
%   \end{equation}
%   It follows inductively that
%   \begin{equation}
%     \|\boldsymbol{h}_l^1-\boldsymbol{h}_l^2\|
%     \leq
%     G_l
%     \|\boldsymbol{\theta}_1-\boldsymbol{\theta}_2\|.
%   \end{equation}
%   Finally, the output layer satisfies
%   \begin{equation}
%     \left\|
%       \Phi(\boldsymbol{x};\boldsymbol{\theta}_1)
%       -
%       \Phi(\boldsymbol{x};\boldsymbol{\theta}_2)
%     \right\|
%     \leq
%     \left(
%       \overline{W}_{L+1}G_L
%       +
%       \overline{h}_L
%       +
%       1
%     \right)
%     \|\boldsymbol{\theta}_1-\boldsymbol{\theta}_2\|.
%   \end{equation}
%   Therefore, the result follows by selecting
%   \begin{equation}
%     L_\Phi(\boldsymbol{x})
%     :=
%     \overline{W}_{L+1}G_L
%     +
%     \overline{h}_L
%     +
%     1.
%   \end{equation}

%   For the componentwise ReLU activation function,
%   \begin{equation}
%     \|\operatorname{ReLU}(\boldsymbol{a})
%       -\operatorname{ReLU}(\boldsymbol{b})\|
%     \leq
%     \|\boldsymbol{a}-\boldsymbol{b}\|,
%   \end{equation}
%   and hence $L_{\phi_l}=1$, which completes the proof.
% \end{proof}
